% Apresentação Beamer sobre uso de álcool e outras drogas entre adolescentes
\documentclass[aspectratio=169]{beamer}
\usetheme{Madrid}
\usepackage[utf8]{inputenc}
\usepackage[brazil]{babel}
\usepackage{graphicx}
\usepackage{amsmath}
\usepackage{booktabs}
\usepackage{hyperref}
\usepackage{tikz}
\setbeamertemplate{navigation symbols}{}

\title{Uso de Álcool e Outras Drogas entre Adolescentes}
\subtitle{Causas, Consequências e Prevenção}
\author{}
\date{}

\begin{document}

\begin{frame}
  \titlepage
\end{frame}

\begin{frame}{Pergunta norteadora}
  \begin{center}
    \Large{\textbf{Por que as pessoas fazem uso de álcool e outras drogas?}}
  \end{center}
\end{frame}

\begin{frame}{Principais motivos para o uso}
  \begin{block}{Motivos sociais}
    \begin{itemize}
      \item Pressão dos amigos, desejo de aceitação.
      \item Festa, celebração e rituais de grupo.
    \end{itemize}
  \end{block}

  \begin{block}{Motivos psicológicos}
    \begin{itemize}
      \item Alívio do stress, ansiedade ou tristeza.
      \item Curiosidade e busca por novas sensações.
    \end{itemize}
  \end{block}
\end{frame}

\begin{frame}{Outros fatores que influenciam}
  \begin{itemize}
    \item Disponibilidade e facilidade de acesso (em casa ou na comunidade).
    \item Modelos familiares (pais ou irmãos que consomem).
    \item Conteúdo divulgado nas redes sociais e na mídia.
    \item Fatores socioeconômicos e falta de atividades de lazer.
  \end{itemize}
\end{frame}

\begin{frame}{Consequências do consumo na adolescência}
  \begin{itemize}
    \item Risco para o desenvolvimento cerebral (memória e atenção).
    \item Problemas escolares: faltas, queda no rendimento.
    \item Risco de acidentes, comportamentos de risco e conflitos familiares.
    \item Possibilidade de dependência e problemas de saúde a longo prazo.
  \end{itemize}
\end{frame}

\begin{frame}{Sinais de que alguém pode estar com problema}
  \begin{itemize}
    \item Mudanças no comportamento: isolamento, irritabilidade, mentiras.
    \item Queda do desempenho escolar e mudanças no sono/alimentação.
    \item Cheiros estranhos, objetos suspeitos, faltar a aulas com frequência.
  \end{itemize}
\end{frame}

\begin{frame}{Como ajudar (semplicidade e respeito)}
  \begin{enumerate}
    \item Conversa calma, sem julgamentos. Ouvir primeiro.
    \item Oferecer apoio e informar sobre riscos (fatos simples).
    \item Procurar a família ou profissionais da escola (psicólogo, orientador).
    \item Se houver risco imediato, buscar ajuda profissional/serviços de saúde.
  \end{enumerate}
\end{frame}

\begin{frame}{Prevenção nas escolas e na comunidade}
  \begin{itemize}
    \item Educação sobre drogas com linguagem clara e baseada em evidências.
    \item Atividades alternativas: esportes, teatro, oficinas e projetos.
    \item Envolver família e professores em programas preventivos.
    \item Políticas escolares claras sobre uso e apoio a quem precisa.
  \end{itemize}
\end{frame}

\end{document}

